\documentclass[12pt,a4paper,twoside]{ctexart}
\usepackage{amsmath,booktabs,geometry,hyperref,needspace,setspace}
\geometry{top=2.5cm,bottom=2.5cm,left=2.7cm,right=2.7cm}
\onehalfspacing
\hypersetup{colorlinks=true,linkcolor=blue,urlcolor=blue}

\title{A bilingual framework for reproducible field studies\\
\zihao{-3}可复现现场研究的双语框架}
\author{Alex Chen / 陈明}
\date{}

\begin{document}
\maketitle

\begin{abstract}
This article presents a compact structure for recording field observations,
quality checks, and interpretation in one bilingual manuscript. The example
keeps the two languages aligned by using the same section sequence and shared
symbols.

\textbf{关键词：}可复现性；现场研究；数据质量
\end{abstract}

\section{English overview}
Field studies often combine observations made by different people and devices.
A document that records the sampling context and the quality decisions beside
the results is easier to audit and update. The following normalized measure is
useful when two instruments have different scales:
\begin{equation}
  y_i = \frac{x_i-x_{\min}}{x_{\max}-x_{\min}}.
\end{equation}

\section{中文概述}
现场研究通常需要合并来自不同人员和设备的观测结果。将采样环境、质量
判断和最终结果放在同一份文档中，可以帮助读者理解分析过程并复核关键
决策。归一化公式中的符号与英文部分保持一致，便于双语读者对照。

\Needspace{10\baselineskip}
\section{Results / 结果}
\begin{table}[ht]
\centering
\caption{Illustrative observations / 示例观测结果}
\begin{tabular}{lrr}
\toprule
Site / 地点 & Samples / 样本 & Valid / 有效(\%)\\
\midrule
River / 河流 & 64 & 94\\
Forest / 森林 & 58 & 90\\
Town / 城镇 & 72 & 96\\
\bottomrule
\end{tabular}
\end{table}

\section{Discussion / 讨论}
The table is intentionally synthetic. Replace it with measurements and explain
the inclusion criteria before using the document as a research record.

表格中的数据仅用于演示排版。正式使用时，应替换为真实测量值，并说明
数据筛选、仪器校准和统计分析的具体方法。

\section{Conclusion / 结论}
Keeping equivalent headings and notation in both languages reduces translation
ambiguity. The starter can be extended with figures, appendices, and a longer
bibliography while retaining the same bilingual workflow.

\begin{thebibliography}{9}
\bibitem{fair} M. D. Wilkinson et al., ``The FAIR Guiding Principles for
scientific data management and stewardship,'' \emph{Scientific Data}, 2016.
\bibitem{latex} L. Lamport, \emph{LaTeX: A Document Preparation System}, 2nd ed.,
Addison--Wesley, 1994.
\end{thebibliography}
\end{document}
